% Chapter introduction: why the topic arose, who first asked, what it answers.
% STATUS: written 2026-11-05 (sources read: Suh 1969 pp. 255--258; Cheung 2001
% pp. 117--127; Cheung 2003 pp. 299--310; Coelho--Milies 1993 pp. 263--267;
% Benkovic 2014 Section 1; Nowicki 1983 pp. 281--297; Alizadeh 2009 pp. 574--578;
% Li--van Wyk--Wei 2012 arXiv version.)
Chapter~\ref{ch:functional-identities} settled the commuting-map and
functional-identity questions on prime and semiprime rings by means of
$d$-free sets \cite[p.~21]{BresarMartindaleChebotar2007FI}. A triangular
algebra is a standard example of a ring that is not $d$-free, so that method
does not reach it \cite[Section~1]{Benkovic2014}. This chapter takes up the
same questions for algebras built from matrices, where the answers come from
explicit computations on the matrix blocks.

The literature of the chapter opens with the 1969 paper of Suh on the Zorn vector
matrix \cite[p.~255]{Suh1969}. Suh replaces the three-dimensional vector
algebra inside the Zorn matrix by a finite-dimensional Lie algebra carrying a
symmetric associative bilinear form \cite[p.~255]{Suh1969}. He then determines
every alternative algebra obtained in this way \cite[p.~255]{Suh1969}. If the
form is nondegenerate the algebra is the split Cayley algebra or a quaternion
algebra \cite[Theorem~1, p.~257]{Suh1969}. The paper shows how a matrix shape
turns a structure problem about the whole algebra into a problem about its
entries.

The derivation question for matrix-shaped algebras has its own line of papers.
For upper triangular matrices over a ring with identity and no further
restriction, Coelho and Milies describe every derivation
\cite[p.~263]{CoelhoMilies1993}. For the subring of $M_n(R)$ cut out by a
reflexive and transitive relation on the index set, Nowicki describes every
derivation and decides when all of them are inner
\cite[Theorem~3.7, p.~287]{Nowicki1983}. For triangular algebras in the abstract
sense, Cheung introduces the algebra and finds all of its commuting maps
\cite[p.~119]{Cheung2001}, and later determines all of its Lie derivations
\cite[Theorem~11, p.~307]{Cheung2003}. The two questions have the same shape: a
map is described block by block, and the part that is not a derivation lands in
the centre.

Section~\ref{sec:matrix-triangular:matrix} reports full matrix algebras and the
special subrings of $M_n(R)$, from the Zorn matrix of Suh through Nowicki
relation-indexed subrings to the Jordan derivations of Alizadeh
\cite[Theorem~3.1, p.~575]{Alizadeh2009}.
Section~\ref{sec:matrix-triangular:triangular} treats triangular algebras, where
the derivation, commuting-map and Lie derivation questions are all settled
\cite[Theorem~11, p.~307]{Cheung2003}.
Section~\ref{sec:matrix-triangular:generalized-matrix} passes to generalized
matrix algebras, which add a second bimodule and two bilinear maps, and reports
the Jordan-derivation theorem of Li, van Wyk and Wei
\cite[Theorem~3.11, p.~13]{LiVanWykWei2012}.
Section~\ref{sec:matrix-triangular:nest-algebraic} closes with the algebraic
side of nest and reflexive algebras, where a triangular algebra appears again as the
compression of a nest algebra at a projection
\cite[Proposition~19, p.~308]{Cheung2003}.

\hbfigure{10-matrix-triangular/triangular}{A triangular algebra: two algebras joined by a bimodule. Derivations split into parts on $A$, on $B$ and on $M$.}{fig:triangular:block}
